%2multibyte Version: 5.50.0.2953 CodePage: 1252

\documentclass{article}
%%%%%%%%%%%%%%%%%%%%%%%%%%%%%%%%%%%%%%%%%%%%%%%%%%%%%%%%%%%%%%%%%%%%%%%%%%%%%%%%%%%%%%%%%%%%%%%%%%%%%%%%%%%%%%%%%%%%%%%%%%%%%%%%%%%%%%%%%%%%%%%%%%%%%%%%%%%%%%%%%%%%%%%%%%%%%%%%%%%%%%%%%%%%%%%%%%%%%%%%%%%%%%%%%%%%%%%%%%%%%%%%%%%%%%%%%%%%%%%%%%%%%%%%%%%%
\usepackage{amsfonts}
\usepackage{amsmath}
\usepackage{adjustbox}
\usepackage[table,xcdraw]{xcolor}
\usepackage{fancyhdr}
\usepackage{xcolor}

\renewcommand\fbox{\fcolorbox{red}{white}}

\usepackage{geometry}
 \geometry{
 a4paper,
 total={170mm,257mm},
 left=20mm,
 top=20mm,
 }
\usepackage{enumitem}


\setcounter{MaxMatrixCols}{10}

\newtheorem{theorem}{Theorem}
\newtheorem{acknowledgement}[theorem]{Acknowledgement}
\newtheorem{algorithm}[theorem]{Algorithm}
\newtheorem{axiom}[theorem]{Axiom}
\newtheorem{case}[theorem]{Case}
\newtheorem{claim}[theorem]{Claim}
\newtheorem{conclusion}[theorem]{Conclusion}
\newtheorem{condition}[theorem]{Condition}
\newtheorem{conjecture}[theorem]{Conjecture}
\newtheorem{corollary}[theorem]{Corollary}
\newtheorem{criterion}[theorem]{Criterion}
\newtheorem{definition}[theorem]{Definition}
\newtheorem{example}[theorem]{Example}
\newtheorem{exercise}[theorem]{Exercise}
\newtheorem{lemma}[theorem]{Lemma}
\newtheorem{notation}[theorem]{Notation}
\newtheorem{problem}[theorem]{Problem}
\newtheorem{proposition}[theorem]{Proposition}
\newtheorem{remark}[theorem]{Remark}
\newtheorem{solution}[theorem]{Solution}
\newtheorem{summary}[theorem]{Summary}
\newenvironment{proof}[1][Proof]{\noindent\textbf{#1.} }{\ \rule{0.5em}{0.5em}}

\begin{document}
\pagestyle{fancy}
\fancyhf{} % sets both header and footer to nothing
\renewcommand{\headrulewidth}{0pt}
\fancyhead{}
\fancyfoot{}
\fancyfoot{}
\fancyfoot[R]{\thepage}
\fancyfoot[L]{}

\thispagestyle{empty}

\begin{center}
{\Large ------ Econometrics ------ \\[1em] \large{Mid-Term Exam} \\[1em] \small{November, 11th 2024}}\bigskip 
\end{center}

\textbf{Each question is worth 2 points.}\\[1em]

\begin{enumerate}
\item State and explain the Gauss-Markov Theorem and its conditions.

\item Describe the important characteristics of the variance of the conditional distribution
of the error term in a linear regression? What are the implications for OLS estimation?

\item What is a control variable, and how does it differ from a variable of interest? Looking at the following table, for what factors are the control variables controlling? Why? Do coefficients on control variables measure causal effects? Explain.

\begin{center}
    \includegraphics[scale=.8]{images/Capture d’écran 2024-11-07 à 16.17.04.png}
\end{center}

\item Consider the following regression model: $Y_i = \beta_0 + \beta_1 X_i + u_i$. Given the definition of $\hat{\beta}_1$, show that  $\hat{\beta}_1 = r_{XY}\frac{s_Y}{s_X}$, where $r_{XY}$ is the sample correlation between $X$ and $Y$ and ${s_X}$ and ${s_Y}$ are the sample standard deviations of $X$ and $Y$. 

\item Consider the following regression model: $Y_i = \beta_0 + \beta_1 X_i + u_i$. Show that the OLS estimator $\hat{\beta}_1$ is an unbiased estimator of $\beta_1$. 

\end{enumerate}\vspace{.25cm}

\end{document}
